\section{Resonancia}
En cuanto a la parte de medir la tension de caida sobre la Resistencia $R$ y a la salida del generador $V_s$, teniendo en cuenta el rango de frecuencias $7,6\,\text{Khz} \leq f_0 \leq 15,2\,\text{Khz}$

\begin{table}[H]
\centering
\caption{Valores obtenidos para la resonancia}
\label{tab:parte2_tp4_traspuesta}
\begin{tabular}{lccccc}
\toprule
\textbf{} & $f_{0}\,[\text{kHz}]$ & $f_{c1}\,[\text{kHz}]$ & $f_{c2}\,[\text{kHz}]$ & $B\,[\text{kHz}]$ & $Q$ \\
\midrule
\textbf{Exp.} & 10.8 & 9.68 & 12.14 & 2.46 & 4.34 \\
\textbf{Teo.} & 12,3   & 11,3   & 13,3    & 2   & 4.39 \\
\bottomrule
\end{tabular}
\end{table}
Valores utilizados para esta parte del experimento: \(L = 76.08\,\text{mH}\), \(C = 2.2\,\text{nF}\) y \(R = 1.1\,\text{k}\Omega\).
\\
\\
PREGUNTAR SOBRE DIFERENCIA ENTRE EXPERIMENTAL Y TEORICO, NO DEJAR ASI!!!!


\newpage